\documentclass{article}[14pt]
\usepackage[T2A]{fontenc}
\usepackage{geometry}
\usepackage{fancyhdr}
\usepackage{tocloft}
\usepackage{hyperref}
\usepackage{titlesec}
\usepackage{graphicx}
\usepackage{caption}
\usepackage{wrapfig}
\usepackage[english, russian]{babel}
\geometry{
	a4paper,
	top=20mm, 
	right=25mm, 
	bottom=25mm, 
	left=25mm
}
\usepackage{amssymb}
\usepackage{fancyvrb}
\usepackage{fvextra}
\usepackage{ dsfont }
\usepackage{amsmath}
\usepackage{listings}
\usepackage{xcolor}
\usepackage{listingsutf8}
\usepackage{pdfpages}
\usepackage{lipsum}
\usepackage{makecell}
\usepackage{float}
\usepackage{longtable}
\usepackage{tabularx}


\lstset{ 
	basicstyle=\ttfamily\small,
	keywordstyle=\color{blue},
	inputencoding=utf8,              
	extendedchars=true,
	stringstyle=\color{red}
	commentstyle=\color{green},
	showspaces=false,         
	showstringspaces=false,  
	numbers=left,              
	numberstyle=\tiny\color{gray},
	breaklines=true,          
	frame=single,           
	captionpos=b,             
	tabsize=4,
	escapeinside={(*@}{@*)}, % для работы с нестандартными текстами
	literate={ё}{{\"o}}1 {Ё}{{\"O}}1 {«}{{\guillemotleft}}1 {»}{{\guillemotright}}1
	{А}{{\CYRA}}1 {Б}{{\CYRB}}1 {В}{{\CYRV}}1 {Г}{{\CYRG}}1 {Д}{{\CYRD}}1 
	{Е}{{\CYRE}}1 {Ж}{{\CYRZH}}1 {З}{{\CYRZ}}1 {И}{{\CYRI}}1 {Й}{{\CYRISHRT}}1 
	{К}{{\CYRK}}1 {Л}{{\CYRL}}1 {М}{{\CYRM}}1 {Н}{{\CYRN}}1 {О}{{\CYRO}}1 
	{П}{{\CYRP}}1 {Р}{{\CYRR}}1 {С}{{\CYRS}}1 {Т}{{\CYRT}}1 {У}{{\CYRU}}1 
	{Ф}{{\CYRF}}1 {Х}{{\CYRH}}1 {Ц}{{\CYRC}}1 {Ч}{{\CYRCH}}1 {Ш}{{\CYRSH}}1 
	{Щ}{{\CYRSHCH}}1 {Ъ}{{\CYRHRDSN}}1 {Ы}{{\CYRERY}}1 {Ь}{{\CYRSFTSN}}1 
	{Э}{{\CYREREV}}1 {Ю}{{\CYRYU}}1 {Я}{{\CYRYA}}1 
	{а}{{\cyra}}1 {б}{{\cyrb}}1 {в}{{\cyrv}}1 {г}{{\cyrg}}1 {д}{{\cyrd}}1 
	{е}{{\cyre}}1 {ж}{{\cyrzh}}1 {з}{{\cyrz}}1 {и}{{\cyri}}1 {й}{{\cyrishrt}}1 
	{к}{{\cyrk}}1 {л}{{\cyrl}}1 {м}{{\cyrm}}1 {н}{{\cyrn}}1 {о}{{\cyro}}1 
	{п}{{\cyrp}}1 {р}{{\cyrr}}1 {с}{{\cyrs}}1 {т}{{\cyrt}}1 {у}{{\cyru}}1 
	{ф}{{\cyrf}}1 {х}{{\cyrh}}1 {ц}{{\cyrc}}1 {ч}{{\cyrch}}1 {ш}{{\cyrsh}}1 
	{щ}{{\cyrshch}}1 {ъ}{{\cyrhrdsn}}1 {ы}{{\cyrery}}1 {ь}{{\cyrsftsn}}1 
	{э}{{\cyrerev}}1 {ю}{{\cyryu}}1 {я}{{\cyrya}}1               
}
\begin{document}
	
\thispagestyle{empty}
\begin{center}
	\Large {\bfseries МИНИСТЕРСТВО НАУКИ И ВЫСШЕГО ОБРАЗОВАНИЯ РОССИЙСКОЙ ФЕДЕРАЦИИ}
	\\
	\vspace{3mm}
	Федеральное государственное автономное образовательное учреждение высшего образования <<Санкт-Петербургский политехнический университет Петра Великого>>
	\\
	\vspace{3mm}
	Институт компьютерных наук и кибербезопасности
	\\
	\vspace{3mm}
	Высшая школа технологий искусственного интеллекта
	\\
	\vspace{3mm}
	Направление: 02.03.01 Математика и компьютерные науки
	\\
	\vspace{2cm}
	\large{
		Методы тестирования программного обеспечения}
	\\
	\LARGE{
		Отчет по лабораторной работе №3\\
		<<Статическое тестирование>>}
\end{center}

\vspace{5.5cm}
\begin{flushleft}
	\Large
	Студент
	\\
	группы 5130201/30101
	\\
	\vspace{2cm}
	\Large
	Преподаватель
	\\ 
\end{flushleft}
\vspace{-4.6cm}
\begin{flushright}
	\Large
	\line(1,0){2cm}
	\hspace{0.4cm}
	Шишмаков В.О.
	\\
	\vspace{2,7cm}
	\hspace{2cm}
	\Large
	\line(1,0){2cm}
	\hspace{0.4cm}
	Курочкин М. А.
	\hspace{-0.1cm}
	
\end{flushright}

\vspace{1cm}
\hspace{9.8cm}
{
	\Large
	\vspace{0.8cm}
	<<\underline{\hspace{1cm}}>>\underline{\hspace{2.5cm}} 2026 г.
}
\hfill \break
\hfill \break
\hfill \break
\hfill \break
\begin{center} \large{Санкт-Петербург, 2026} \end{center}
\pagebreak

\titleformat{\section}
{\Huge\bfseries}
{\thesection\;}
{0.65cm}
{}
\titleformat{\subsection}
{\huge\bfseries}
{\thesubsection\;}
{0.65cm}
{}
\titleformat{\subsubsection}
{\Large\bfseries}
{\thesubsubsection\;}
{0.65cm}
{}
\titlespacing*{\section}
{0em}
{1.5ex plus 1ex minus 0.2ex}
{1.5ex plus 1ex minus 0.2ex}	

\renewcommand*\contentsname{\huge{Содержание}}
\Large{
	\tableofcontents}
\pagebreak


\section*{Введение}
\addcontentsline{toc}{section}{Введение}

Долгое время большинство программистов полагало, что программы создаются только для выполнения на машинах и не рассчитаны на чтение людьми, 
а единственным способом их проверки служит запуск кода на компьютере. Эта ситуация начала меняться в начале 1970-х годов благодаря тем 
разработчикам, которые первыми поняли ценность просмотра кода как части процесса тестирования и отладки.

Полезность вычитки кода определяется несколькими факторами: размер и сложность приложения, количество членов команды разработчиков, 
временные рамки проекта (например, мягкие или жёсткие), а также, безусловно, квалификация и профессиональная культура команды программистов.

Статический анализ кода представляет собой анализ исходного кода программного обеспечения, выполняемый без фактического запуска 
исследуемых программ.

Спрсcheck --- это инструмент статического анализа для языков C и C++, который предназначен для выявления проблем, не обнаруживаемых 
компиляторами.

\pagebreak

\section{Постановка задачи}

\noindent \textbf{Дано:} программный код на языке программирования \texttt{C++} (Построение словаря на основе хэш-функции и красно-черного дерева. Кодирование данных).

\noindent \textbf{Требуется:} изучить метод статического анализа кода, провести статический анализ кода при помощи \texttt{CppCheck} версии 2.20.0.

\pagebreak

\section{Описание тестируемой программы}

Название программы: нахождение кратчайшего пути между двумя вершинами графа методом обхода в ширину (BFS).
 
\textbf{Дано:}
 
\begin{itemize}
    \item \texttt{n} --- количество вершин графа;
    \item \texttt{oriented} --- тип графа: ориентированный или неориентированный;
    \item \texttt{edges} --- список рёбер графа, задаваемых пользователем вручную (пары вершин);
    \item \texttt{start} --- начальная вершина пути;
    \item \texttt{end} --- конечная вершина пути.
\end{itemize}
 
\textbf{Требуется:} найти кратчайший путь (по числу рёбер) от вершины \texttt{start} до вершины \texttt{end} в заданном графе методом обхода в ширину и вывести последовательность вершин этого пути.
 
\textbf{Ограничения} --- все параметры целые числа:
 
\begin{itemize}
    \item количество вершин $n \in [2;\, 10000]$
    \item тип графа $\mathtt{oriented} \in \{y,\, n\}$
    \item начальная вершина $\mathtt{start} \in [0;\, n-1]$
    \item конечная вершина $\mathtt{end} \in [0;\, n-1],\quad \mathtt{end} \neq \mathtt{start}$
    \item вершины рёбер $u, v \in [0;\, n-1],\quad u \neq v$
\end{itemize}
 
\subsection*{Спецификация}
 
\begin{longtable}{|>{\raggedright\arraybackslash}p{3.2cm}
                  |>{\raggedright\arraybackslash}p{4.5cm}
                  |>{\raggedright\arraybackslash}p{7.3cm}|}
\hline
\textbf{Входные данные} & \textbf{Выходные данные} & \textbf{Реакция программы} \\
\hline
\endfirsthead
\hline
\textbf{Входные данные} & \textbf{Выходные данные} & \textbf{Реакция программы} \\
\hline
\endhead
 
$n < 2$ или $n > 10000$ &
Сообщение: \textit{<<Incorrect input! Enter a number from 2 to 10000.>>} &
Повторный запрос ввода от пользователя до тех пор, пока не будет введено корректное значение. \\
\hline
 
$\mathtt{oriented} \notin \{y, n, Y, N\}$ &
Сообщение: \textit{<<Incorrect input! Enter y or n.>>} &
Повторный запрос ввода от пользователя до тех пор, пока не будет введено корректное значение. \\
\hline
 
Ребро $u \notin [0;\, n{-}1]$ или $v \notin [0;\, n{-}1]$ &
Сообщение: \textit{<<Incorrect input!>>} &
Введённое ребро отклоняется; программа запрашивает следующее ребро. \\
\hline
 
Ребро $u = v$ &
Сообщение: \textit{<<Self-loops are not allowed!>>} &
Введённое ребро отклоняется; программа запрашивает следующее ребро. \\
\hline
 
Ввод $-1$ в поле <<From>> или <<To>> &
Ввод рёбер завершается &
Программа переходит к запросу начальной и конечной вершин. \\
\hline
 
$\mathtt{start} \notin [0;\, n{-}1]$ &
Сообщение: \textit{<<Incorrect input! Enter a number from 0 to $n{-}1$.>>} &
Повторный запрос ввода до тех пор, пока не будет введено корректное значение. \\
\hline
 
$\mathtt{end} \notin [0;\, n{-}1]$ или $\mathtt{end} = \mathtt{start}$ &
Сообщение: \textit{<<Incorrect input! Enter a number from 0 to $n{-}1$ (not start).>>} &
Повторный запрос ввода до тех пор, пока не будет введено корректное значение. \\
\hline
 
Все параметры корректны; путь от \texttt{start} до \texttt{end} существует &
Сообщение вида: \textit{<<Shortest path ($k$ edges): $v_0 \to v_1 \to \dots \to v_k$>>} &
Программа выполняет обход в ширину (BFS) по матрице смежности. Выводит последовательность вершин кратчайшего пути и число пройденных рёбер $k$. Завершает работу. \\
\hline
 
Все параметры корректны; путь от \texttt{start} до \texttt{end} \textbf{не существует} &
Сообщение: \textit{<<No path from start to end.>>} &
Программа сообщает об отсутствии пути и завершает работу. \\
\hline
 
\end{longtable}

\pagebreak

\section{Статический анализ кода}

Статический анализ кода проводится без фактического запуска программы. Он исследует исходный код для обнаружения возможных ошибок, 
уязвимостей и отклонений от принятых стандартов кодирования. Благодаря этому виду анализа разработчики могут находить проблемы на ранних 
этапах создания программы, что позволяет существенно снизить затраты на их исправление в дальнейшем.

Помимо прочего, статический анализ может включать проверку стиля оформления кода, что способствует поддержанию единообразия в проекте. 
Это особенно важно в командах, где над одним проектом трудится несколько разработчиков. Инструменты статического анализа можно настраивать 
на соответствие кода определённым стандартам — это помогает избегать типичных ошибок и повышает удобочитаемость кода.

Многие подобные инструменты можно встраивать в процесс сборки, что позволяет автоматизировать проверку кода. Такой подход обеспечивает 
непрерывный контроль качества и даёт разработчикам возможность больше сосредоточиться на написании кода, а не на его ручной проверке.

\subsection*{Виды статического тестирования}

\textbf{Статический анализ кода (Static Code Analysis)} — этот вид статического тестирования подразумевает исследование исходного кода. 
Его главная цель — выявление потенциальных проблем, некорректных практик, структурных аномалий и нарушений стандартов кодирования. 
Такие инструменты, как Lint, Pylint и ESLint, помогают автоматизировать этот процесс.

\textbf{Обзоры кода (Code Reviews)} — данный вид статического тестирования заключается в анализе кода членами команды разработки или 
приглашёнными экспертами. Обзоры кода позволяют находить проблемы и несоответствия стандартам, а также способствуют обмену знаниями и 
опытом между участниками команды.

\textbf{Анализ архитектуры (Architecture Analysis)} — при этом виде тестирования анализируется архитектура программного обеспечения: её 
структура, связи между компонентами и соблюдение архитектурных принципов. Это помогает выявить проблемы, связанные с проектированием системы.

\pagebreak

\section{Описание приложения и среды разработки}

CppCheck представляет собой мощное средство статического анализа кода, написанного на языках C и C++. Оно позволяет обнаруживать 
разнообразные дефекты, уязвимости и места для оптимизации без необходимости компиляции или запуска программы. Проверка кода может 
выполняться как через командную строку, так и с помощью графического интерфейса.

\subsection*{Ключевые возможности CppCheck}

Анализатор классифицирует найденные проблемы по уровням серьёзности:

\begin{itemize}
\item \textbf{Error} — критические ошибки, способные вызвать сбой программы.
\item \textbf{Warning} — потенциальные проблемы, которые могут привести к некорректному поведению.
\item \textbf{Style} — нарушения стиля оформления, не влияющие на исполнение, но ухудшающие читаемость кода.
\item \textbf{Performance} — рекомендации по повышению производительности.
\item \textbf{Portability} — вопросы, связанные с переносимостью кода на разные платформы.
\item \textbf{Information} — информационные сообщения.
\end{itemize}

Для проверки кода по всем категориям можно использовать следующую команду:

\begin{verbatim}
cppcheck --enable=all --inconclusive исходные_файлы
\end{verbatim}

Чтобы включить только определённые типы проблем, применяются соответствующие ключи:

\begin{itemize}
\item \texttt{--enable=style} — проверка стиля кода.
\item \texttt{--enable=performance} — проверка производительности.
\item \texttt{--enable=portability} — проверка переносимости.
\item \texttt{--enable=information} — информационные сообщения.
\item \texttt{--enable=unusedFunction} — поиск неиспользуемых функций.
\item \texttt{--enable=warning} — вывод предупреждений.
\end{itemize}

\subsection*{Форматы вывода}

CppCheck поддерживает два формата представления результатов анализа:
\begin{itemize}
\item текстовый;
\item XML.
\end{itemize}

\subsection*{Правила проверки}

Анализатор реализует 167 правил, сгруппированных по следующим категориям:

\begin{itemize}
\item 64-bit portability
\item Assert
\item Auto Variables
\item Boolean
\item Boost usage
\item Bounds checking
\item Check function usage
\item Class
\item Condition
\item Exception Safety
\item IO using format string
\item Leaks (auto variables)
\item Memory leaks (address not taken)
\item Memory leaks (class variables)
\item Memory leaks (function variables)
\item Memory leaks (struct members)
\item Null pointer
\item Other
\item STL usage
\item Sizeof
\item String
\item Type
\item Uninitialized variables
\item Unused functions
\item UnusedVar
\item Using postfix operators
\item Vaarg
\end{itemize}

\subsection*{Разбор правил группы «STL usage»}

В группе «STL usage» CppCheck предоставляет 13 проверок на некорректное использование функций стандартной библиотеки шаблонов. 
Рассмотрим два из них.

\subsubsection*{Правило \texttt{dereferencingAnErasedIterator}}
\begin{itemize}
\item Уровень серьёзности: Error
\item Категория: STL usage
\end{itemize}
В STL действия с контейнерами обычно выполняются через итераторы (специальные объекты-посредники, предоставляющие доступ к 
элементам и операции над ними). Однако итераторы не всегда устойчивы: если после создания итератора контейнер был изменён, 
итератор может стать недействительным, и его использование недопустимо. Поведение зависит от типа контейнера. Например, для 
\texttt{std::vector} после вставки нового элемента все итераторы становятся невалидными, а для \texttt{std::list} они остаются 
корректными. Это правило выявляет проблемы, связанные с использованием итераторов после модификации контейнера.

\subsubsection*{Правило \texttt{usingAutoPointer}}
\begin{itemize}
\item Уровень серьёзности: Warning
\item Категория: STL usage
\end{itemize}
STL предоставляет несколько классов для автоматического управления ресурсами: \texttt{std::unique\_ptr} (уникальное владение) и 
\texttt{std::shared\_ptr} (разделяемое владение). Оба появились в стандарте C++11. Более старый стандарт C++03 включал только 
\texttt{std::auto\_ptr} (также реализующий уникальное владение), но этот класс признан устаревшим. Его использование может 
приводить к проблемам с управлением ресурсами, поэтому считается нежелательным. Данное правило отслеживает применение 
\texttt{std::auto\_ptr}.

\subsection*{Разбор правил группы «Bounds checking»}

Эта группа проверяет корректность работы с массивами и строками. Рассмотрим два правила.

\subsubsection*{Правило \texttt{arrayIndexOutOfBounds}}
\begin{itemize}
\item Уровень серьёзности: Error
\item Категория: Bounds checking
\end{itemize}
Правило действует для статических массивов, а также для динамически выделенных массивов, размер которых известен на этапе компиляции. 
Оно проверяет, что при обращении к массиву не происходит выхода за его границы.

\subsubsection*{Правило \texttt{negativeMemoryAllocationSize}}
\begin{itemize}
\item Уровень серьёзности: Error
\item Категория: Bounds checking
\end{itemize}
Выделение динамического или статического массива с отрицательным размером является ошибочным. Компилятор может обнаружить такую проблему 
только если отрицательное значение задано явно в коде. Если же размер передаётся через переменную, компилятор обычно не может выявить ошибку. 
Однако данное правило CppCheck способно определить проблему в тех случаях, когда значение переменной можно вычислить статически.

\pagebreak

\section{Запуск CppCheck на MacOS}

Работа с CppCheck на MacOS происходит в терминале. Для того, чтобы установить и работать с программой требуется следовать следующим шагам:

\begin{enumerate}
	\item Откройте терминал.
	\item Для установки воспользуйтесь командой:
	
	\begin{verbatim}
		brew install cppcheck
	\end{verbatim}

	\item Выполните анализ программы при помощи команды:
	
	\begin{verbatim}
		cppcheck --enable=all --inconclusive путь_к_исходным_файлам
	\end{verbatim}

\end{enumerate}

Для того просмотра ошибок определенного типа используйте следующие параметры:

\begin{itemize}
\item \texttt{enable=style} — анализ соответствия стилевым нормам кода;
\item \texttt{enable=performance} — оценка производительности;
\item \texttt{enable=portability} — контроль переносимости на различные платформы;
\item \texttt{enable=information} — вывод информационных сообщений;
\item \texttt{enable=unusedFunction} — обнаружение функций, которые нигде не используются;
\item \texttt{enable=warning} — показ предупреждений.
\end{itemize}

\pagebreak

\section{Результаты запуска}

После запуска анализатора Cppcheck с параметрами \texttt{--enable=all --inconclusive}
для файлов \texttt{main.cpp}, \texttt{graph.cpp} и \texttt{graph.h} было обнаружено:
1 ошибка синтаксиса, 1 стилистическое предупреждение и 7 информационных сообщений
о системных заголовочных файлах.
 
Результаты анализа представлены в таблице ниже. В таблице указаны файл и строка,
в которой обнаружена проблема, важность, краткое описание и идентификатор проблемы.

{\scriptsize{
\begin{longtable}{|>{\raggedright\arraybackslash}p{2.2cm}
                  |>{\raggedright\arraybackslash}p{1.3cm}
                  |>{\raggedright\arraybackslash}p{2.4cm}
                  |>{\raggedright\arraybackslash}p{5.5cm}
                  |>{\raggedright\arraybackslash}p{3.0cm}|}
\hline
\textbf{Файл} & \textbf{Строка} & \textbf{Важность} & \textbf{Кратко} & \textbf{Id} \\
\hline
\endfirsthead
\hline
\textbf{Файл} & \textbf{Строка} & \textbf{Важность} & \textbf{Кратко} & \textbf{Id} \\
\hline
\endhead
 
\texttt{graph.h}   & 7  & ошибка         & Code `classmyGraph\{' is invalid C code. & \texttt{syntaxError} \\
\hline
\texttt{graph.cpp} & 14 & стиль          & The function `getNum' is never used. & \texttt{unusedFunction} \\
\hline
\texttt{main.cpp}  & 1  & информация & Include file: \textless iostream\textgreater{} not found. & \texttt{missingIncludeSystem} \\
\hline
\texttt{main.cpp}  & 2  & информация & Include file: \textless string\textgreater{} not found. & \texttt{missingIncludeSystem} \\
\hline
\texttt{main.cpp}  & 3  & информация & Include file: \textless regex\textgreater{} not found. & \texttt{missingIncludeSystem} \\
\hline
\texttt{graph.h}   & 3  & информация & Include file: \textless vector\textgreater{} not found. & \texttt{missingIncludeSystem} \\
\hline
\texttt{graph.h}   & 4  & информация & Include file: \textless iostream\textgreater{} not found. & \texttt{missingIncludeSystem} \\
\hline
\texttt{graph.h}   & 5  & информация & Include file: \textless queue\textgreater{} not found. & \texttt{missingIncludeSystem} \\
\hline
\texttt{main.cpp}  & 0  & информация & Limiting analysis of branches. Use --check-level=exhaustive. & \texttt{normalCheckLevel\-MaxBranches} \\
\hline
 
\end{longtable}}}
 
\subsection*{Описание найденных проблем}
 
\begin{enumerate}
 
\item \textbf{Ошибка синтаксиса:}
 
Сообщение: \texttt{Code 'classmyGraph\{' is invalid C code.}
 
Описание: Cppcheck анализировал файл \texttt{graph.h} как код на языке C, а не C++.
При таком режиме анализа ключевое слово \texttt{class} не распознаётся,
и инструмент воспринимает строку \texttt{class myGraph \{} как единый некорректный
идентификатор \texttt{classmyGraph\{}. Реальной синтаксической ошибки в коде нет.
 
Решение: передавать файл на анализ явно как C++ код, добавив флаг
\texttt{--language=c++}, либо убедиться, что расширение файла
(\texttt{.h}, \texttt{.cpp}) корректно распознаётся анализатором.
 
\item \textbf{Стилистическое предупреждение:}
 
Сообщение: \texttt{The function 'getNum' is never used.}
 
Описание: метод \texttt{getNum()} объявлен и реализован в классе \texttt{myGraph},
однако нигде не вызывается в коде программы. После упрощения программы до
нахождения кратчайшего пути метод утратил своё назначение и стал <<мёртвым кодом>>.
 
Решение: удалить объявление метода \texttt{getNum()} из файла \texttt{graph.h}
и его реализацию из файла \texttt{graph.cpp}.
 
\item \textbf{Информационные сообщения (\texttt{missingIncludeSystem}):}
 
Сообщение (7 случаев): \texttt{Include file: <...> not found.}
 
Описание: Cppcheck не находит стандартные заголовочные файлы
(\texttt{<iostream>}, \texttt{<string>}, \texttt{<regex>}, \texttt{<vector>},
\texttt{<queue>}). Как указывает сам анализатор, стандартные библиотечные
заголовки не обязательно предоставлять для получения корректных результатов
анализа. Это не ошибки в коде, а особенность окружения запуска Cppcheck.
 
Решение: сообщения носят исключительно информационный характер и не требуют
исправления кода. При необходимости их можно подавить флагом
\texttt{--suppress=missingIncludeSystem}.
 
\end{enumerate}

\pagebreak

\section{Исправление недостатков}

Изменения:
 
\begin{enumerate}
 
\item \textbf{Удаление неиспользуемого метода \texttt{getNum()} (\texttt{unusedFunction}):}
метод был объявлен в \texttt{graph.h} и реализован в \texttt{graph.cpp},
однако нигде не вызывался. Метод удалён из обоих файлов.
 
Было в \texttt{graph.h}:
\begin{lstlisting}
    std::vector<int> bfsShortestPath(int start, int end);
    int getNum() const;
};
\end{lstlisting}
 
Стало в \texttt{graph.h}:
\begin{lstlisting}
    std::vector<int> bfsShortestPath(int start, int end);
};
\end{lstlisting}
 
Было в \texttt{graph.cpp}:
\begin{lstlisting}
int myGraph::getNum() const {
    return m_num;
}
\end{lstlisting}
 
Стало в \texttt{graph.cpp}: реализация метода полностью удалена.
 
\item \textbf{Исправление ложной ошибки синтаксиса (\texttt{syntaxError}):}
Cppcheck анализировал \texttt{graph.h} как код на языке C, из-за чего
ключевое слово \texttt{class} не распознавалось и возникала ложная ошибка.
Решение~--- не передавать заголовочный файл явно, так как он подключается
автоматически через \texttt{\#include} в \texttt{.cpp} файлах.
 
Было:
\begin{lstlisting}
cppcheck --enable=all --inconclusive main.cpp graph.cpp graph.h
\end{lstlisting}
 
Стало:
\begin{lstlisting}
cppcheck --enable=all --inconclusive main.cpp graph.cpp
\end{lstlisting}
 
\end{enumerate}
 
После исправления всех проблем Cppcheck не выявляет ни одной ошибки,
стилистического предупреждения или проблемы производительности.
Результат финального запуска анализатора:
 
\begin{lstlisting}
user@MacBook-Air-user fixed-app % cppcheck --enable=all --inconclusive main.cpp graph.cpp
Checking main.cpp ...
main.cpp:1:2: information: Include file: <iostream> not found. [missingIncludeSystem]
main.cpp:2:2: information: Include file: <string> not found. [missingIncludeSystem]
main.cpp:3:2: information: Include file: <regex> not found. [missingIncludeSystem]
graph.h:3:2: information: Include file: <vector> not found. [missingIncludeSystem]
graph.h:4:2: information: Include file: <queue> not found. [missingIncludeSystem]
main.cpp:0:0: information: Limiting analysis of branches. [normalCheckLevelMaxBranches]
1/2 files checked 74% done
Checking graph.cpp ...
2/2 files checked 100% done
nofile:0:0: information: Active checkers: 183/186 [checkersReport]
\end{lstlisting}
 
Итоговая статистика финального сканирования:
 
\begin{longtable}{|>{\raggedright\arraybackslash}p{9cm}
                  |>{\raggedright\arraybackslash}p{3cm}|}
\hline
\textbf{Категория} & \textbf{Количество} \\
\hline
\endfirsthead
\hline
\textbf{Категория} & \textbf{Количество} \\
\hline
\endhead
Ошибки & 0 \\
\hline
Предупреждения & 0 \\
\hline
Стилистические предупреждения & 0 \\
\hline
Предупреждения переносимости & 0 \\
\hline
Проблемы с производительностью & 0 \\
\hline
Информационные сообщения (\texttt{missingIncludeSystem}) & 5 \\
\hline
Активных анализаторов & 183 из 186 \\
\hline
\end{longtable}

\pagebreak

\section*{Заключение}
\addcontentsline{toc}{section}{Заключение}

В ходе выполнения лабораторной работы был проведён статический анализ кода
при помощи Cppcheck, а также описана сама программа. Использование Cppcheck
позволило обнаружить одну ошибку синтаксиса и одно стилистическое предупреждение.
На мой взгляд, Cppcheck~--- это удобный инструмент, которым легко пользоваться
для тестирования программного кода.
 
Ошибка синтаксиса оказалась ложной: анализатор обрабатывал заголовочный файл
\texttt{graph.h} как код на языке C, из-за чего ключевое слово \texttt{class}
не распознавалось корректно. Проблема была устранена изменением команды запуска
анализатора~--- без явной передачи \texttt{graph.h} в аргументах. Стилистическое
предупреждение об неиспользуемом методе \texttt{getNum()} было устранено путём
удаления этого метода из кода, поскольку после упрощения программы он утратил
своё назначение.
 
Важно понимать, что статическое тестирование с помощью Cppcheck позволяет
автоматически выявлять проблемы, которые сложно заметить при ручной проверке.
Опыт, заложенный в инструменты статического анализа, несопоставим с возможностями
одного программиста. Я сделал вывод, что статическое тестирование является
полезным и необходимым этапом разработки, позволяющим повысить качество кода
и устранить потенциальные ошибки на ранней стадии.

\end{document}